\documentclass[11pt]{article}
\usepackage[letterpaper,margin=1in]{geometry}
\usepackage[T1]{fontenc}
\usepackage{amsmath,amssymb,amsthm,mathtools}
\usepackage{lmodern}
\usepackage{microtype}
\usepackage[colorlinks=true,linkcolor=blue,citecolor=blue,urlcolor=blue]{hyperref}
\newtheorem{theorem}{Theorem}
\newtheorem{lemma}[theorem]{Lemma}
\newtheorem{proposition}[theorem]{Proposition}
\newtheorem{corollary}[theorem]{Corollary}
\theoremstyle{definition}
\newtheorem{definition}[theorem]{Definition}
\newtheorem{remark}[theorem]{Remark}
\DeclareMathOperator{\conv}{conv}
\DeclareMathOperator{\aff}{aff}
\DeclareMathOperator{\intr}{int}
\DeclareMathOperator{\relint}{rel\,int}
\DeclareMathOperator{\cl}{cl}
\DeclareMathOperator{\xc}{xc}
\DeclareMathOperator{\proj}{proj}
\newcommand{\R}{\mathbb R}
\newcommand{\Z}{\mathbb Z}
\newcommand{\PM}{P_{\mathrm M}}
\title{Extension Complexity Lower Bounds\\for Mixed-Integer Extended Formulations\\[0.4em]
\large A corrected argument using families of formulations}
% Author names are retained from the supplied 2016 source.
% Historical affiliations and email addresses are omitted pending author review.
\author{Robert Hildebrand \and Robert Weismantel \and Rico Zenklusen}
\date{Revised draft --- October 7, 2026}

\begin{document}
\maketitle

\begin{abstract}
We revisit a geometric approach to lower bounds on the number of integer
variables in mixed-integer linear extended formulations. We prove that every
formulation of the matching polytope of the complete graph on $n$ vertices
with polynomially many inequalities requires
$\Omega(\log n/\log\log n)$ integer variables. The same conclusion holds
for the perfect matching and traveling salesman polytopes. The proof combines
the exponential linear extension complexity of the matching polytope,
lattice-free faces, and flatness. To avoid an invalid interchange of convexification
and integrality in the earlier version, we retain a family of formulations
throughout the argument and control its evolution by a decreasing potential.
This gives a weaker bound than the one claimed in that version.
The substantially stronger results of Cevallos, Weltge, and Zenklusen subsume
these lower bounds. Our purpose is to correct the earlier work and present
an alternative proof using exact faces and the original geometric ingredients.
An appendix identifies the prior error and gives a counterexample to the
construction used there.
\end{abstract}

\section{Introduction}

Mixed-integer linear extended formulations (MILEFs) are a standard framework
for modeling problems in discrete optimization. A natural measure of their
complexity is the number of integer variables needed when the number of linear
inequalities is small. This question is distinct from the number of variables
in a particular familiar formulation: auxiliary continuous variables and
alternative encodings can change the number of integrality conditions
substantially. The traveling salesman problem provides many examples of
different mixed-integer descriptions; see, for instance,
\cite{miller_1960_integer,padberg_1991_analytical,oncan_2009_comparative}.

Our starting point is the exponential lower bound on the linear extension
complexity of the matching polytope due to Rothvo\ss~\cite{rothvoss_2014_matching}.
We combine it with flatness for lattice-free convex sets and the behavior of
matching polytopes under restrictions to faces supported on a subset of the
vertices. These are the central geometric ingredients of the earlier version
of this paper~\cite{HWZ2016}.

\subsection{Purpose of this revision and subsequent work}
The earlier version claimed a lower bound of
$\Omega(\sqrt{n/\log n})$ integer variables for polynomial-size MILEFs of
the matching polytope. Its proof of Lemma~8 incorrectly interchanged taking
a convex hull and imposing the remaining integrality conditions. Consequently,
the proposed elimination of an integer variable did not preserve the projected
mixed-integer hull. The earlier version was subsequently withdrawn.
Appendix~\ref{app:error} gives the precise invalid equality and a concrete
counterexample.

Here we replace that elimination step by an induction over a \emph{family}
of formulations. After choosing a suitable face, we split one member into
finitely many slices and retain those slices separately. We never replace
their union by its linear convex hull while retaining integrality conditions
on averaged variables. This change repairs the argument, at the cost of a
weaker quantitative conclusion.

Cevallos, Weltge, and Zenklusen~\cite{CWZ2018} established the much stronger
lower bound $\Omega(n/\log n)$ for matching and traveling salesman polytopes,
even for formulations of subexponential size. Their work provides a general
framework transferring lower bounds for approximate linear extended formulations
to the mixed-integer setting, and applies to additional families of polyhedra.
In particular, their results subsume the polynomial-size lower bounds proved
here, as well as the stronger asymptotic claim made in the earlier version.
We do not claim an improvement of those results. The purpose of this revision
is to put the earlier geometric approach on a correct footing and to provide
an alternative derivation that uses exact matching faces and exact linear
extension complexity, without invoking their approximation framework.

\subsection{Formulations and the main result}
We use affine maps to state formulations; this makes equations and changes of
integer coordinates convenient to handle.

\begin{definition}
A MILEF of a polytope $P\subseteq\R^d$ is a triple $(Q,\sigma,\pi)$, where
$Q\subseteq\R^p$ is a polyhedron and $\sigma:\R^p\to\R^k$ and
$\pi:\R^p\to\R^d$ are affine maps, such that
\begin{equation}\label{eq:milef}
 P=\conv\bigl(\pi(Q\cap\sigma^{-1}(\Z^k))\bigr).
\end{equation}
Its size is at most $(m,k)$ if $Q$ has a description with at most $m$
linear inequalities, in addition to arbitrary linear equations.
\end{definition}

This convention is equivalent to requiring selected coordinates to be integer:
introduce variables $z=\sigma(y)$ and impose $z\in\Z^k$. Similarly, introducing
$x=\pi(y)$ makes the output map a coordinate projection. These changes add
only equations. We retain equations explicitly rather than eliminating them
in a way that could alter the integrality conditions. Coefficients may be real;
no rationality assumption is used in the proof. The argument is existential
and makes no claim of an efficient algorithm for constructing its intermediate
families.

For a finite vertex set $V$, let $E(V)$ denote the edges of the complete graph
on $V$, and let $\PM(V)\subseteq\R^{E(V)}$ be its matching polytope, the
convex hull of the incidence vectors of all matchings. This includes the empty
matching. In particular, for $|V|\leq1$, the polytope is a singleton in a
zero-dimensional space, not the empty set.

\begin{theorem}\label{thm:main}
Fix a polynomial $p$. There exist constants $c_p>0$ and $n_p$ such that, for
every $n\geq n_p$, every MILEF of the matching polytope of $K_n$ with at most
$p(n)$ linear inequalities has at least
\[
 c_p\frac{\log n}{\log\log n}
\]
integer variables.
\end{theorem}

\begin{corollary}\label{cor:other}
The conclusion of Theorem~\ref{thm:main} also holds for the perfect matching
polytope of $K_n$, with $n$ even, and for the traveling salesman polytope of
$K_n$.
\end{corollary}

\begin{proof}
The matching polytope of $K_t$ is a projection of the perfect matching polytope
of $K_{2t}$: use two copies of the vertex set, match the unmatched vertices
across the copies, and complete using the same matching in each copy.
The standard reduction to the traveling salesman polytope uses a graph with
$O(t)$ vertices; see~\cite{yannakakis_1991_expressing,rothvoss_2014_matching}.
These projections, and restrictions to faces if needed in the reduction,
preserve the number of integer variables and preserve a polynomial bound
on the number of inequalities. Passing to a smaller complete graph by fixing
edges (and, for tours, contracting a fixed path) handles intermediate graph
sizes. Applying Theorem~\ref{thm:main} proves the claim.
\end{proof}

The theorem excludes a constant number of integer variables, and more generally
$o(\log n/\log\log n)$ integer variables, in polynomial-size formulations.
It does not exclude $\Theta(\log n/\log\log n)$ integer variables.
Thus it does not, by itself, rule out the parameter range in which a
$k^{O(k)}$ dependence in fixed-dimension mixed-integer optimization is
polynomial in $n$; compare~\cite{lenstra_integer_1983}.

Related questions include affine TU-dimension~\cite{bader_2016_mixed} and
the size of integer descriptions without additional variables
\cite{kaibel_2015_lower}. Our argument concerns arbitrary extended
formulations and uses only their inequality count and the number of
integrality conditions. Throughout, logarithms are to base $2$, and convex
hulls consist of finite convex combinations; they are not implicitly closed.

\section{Geometric ingredients}\label{sec:ingredients}

We first record the extension-complexity and flatness bounds used below.

\begin{theorem}[Rothvo\ss~\cite{rothvoss_2014_matching}]\label{thm:rothvoss}
There is a constant $c>0$ such that
$\xc(\PM(V))\geq 2^{c|V|}$ whenever $|V|\geq2$.
\end{theorem}

\begin{theorem}[Flatness; see~\cite{Khinchin1948,Banaszczyk1999}]\label{thm:flatness}
There is an absolute constant $C\geq1$ such that every closed,
full-dimensional convex set $K\subseteq\R^s$, with $s\geq1$ and
$\intr(K)\cap\Z^s=\varnothing$, has a nonzero integer direction $v$ with
\[
 \sup_{z\in K}v^Tz-\inf_{z\in K}v^Tz\leq Cs^{3/2}.
\]
\end{theorem}

The statement also allows unbounded $K$. Indeed, intersect $K$ with growing
boxes containing a fixed ball in its interior. Flatness gives directions for
these bounded intersections. Their norms are bounded because the fixed ball
has bounded width in each such direction. One integer direction therefore
occurs infinitely often and has the asserted width for $K$.
Fix $C$ as above and define, for $s\geq1$,
\begin{equation}\label{eq:branching}
 b(s)=\lceil Cs^{3/2}\rceil+1.
\end{equation}

\subsection{Matching faces}

\begin{lemma}\label{lem:matchingface}
Let $W\subseteq V$, and let $\bar F$ be a nonempty face of $\PM(W)$.
Then the face
\[
 F=\{x\in\PM(V):x_{E(W)}\in\bar F\}
\]
projects onto $\PM(V\setminus W)$ under the projection onto $E(V\setminus W)$.
\end{lemma}

\begin{proof}
Restricting any matching on $V$ to $E(V\setminus W)$ gives a matching.
Conversely, choose a matching $\bar M$ on $W$ whose incidence vector lies in
$\bar F$. For any matching $M$ on $V\setminus W$, the union $M\cup\bar M$
is a matching on $V$ whose incidence vector lies in $F$ and projects to the
incidence vector of $M$. Taking convex hulls gives the assertion.
\end{proof}

\subsection{A lattice-free set from a violated valid inequality}
The following form of the lattice-free-face argument is essential: an individual
member of a family need not describe the whole target polytope.

\begin{lemma}\label{lem:latticefree}
Let $Q\subseteq\R^p$ be convex, let $\sigma:\R^p\to\R^s$ be affine with
$s\geq1$, and let $h:\R^p\to\R$ be affine. Suppose that
\[
 h(y)\leq0\quad\text{for all }y\in Q\cap\sigma^{-1}(\Z^s),
 \qquad h(q)>0\quad\text{for some }q\in Q.
\]
Write $H=\{y:h(y)=0\}$. Then
\[
 K=\conv\bigl(\sigma(Q\cap H\cap\sigma^{-1}(\Z^s))\bigr)
\]
satisfies $\intr(K)\cap\Z^s=\varnothing$.
\end{lemma}

\begin{proof}
Suppose $\bar z\in\intr(K)\cap\Z^s$. Choose affinely independent
$z^1,\ldots,z^{s+1}\in K$ such that $\bar z$ is a strict convex combination
$\sum_i\lambda_i z^i$. Since affine maps commute with convex hulls, choose
lifts
\[
 y^i\in\conv(Q\cap H\cap\sigma^{-1}(\Z^s)),
 \qquad \sigma(y^i)=z^i.
\]
Put $y=\sum_i\lambda_i y^i$, $Y=\conv\{y^1,\ldots,y^{s+1}\}$, and
$A=\aff(Y)$. The restriction of $\sigma$ to $A$ is an affine bijection
onto $\R^s$, because the images $z^i$ are affinely independent. In particular,
there is $w\in A$ with $\sigma(w)=\sigma(q)$. Also $A\subseteq H$,
$y\in\relint(Y)$, and $\sigma(y)=\bar z$.

For sufficiently small $\epsilon\in(0,1)$, define
\[
 y_\epsilon=y+\epsilon(q-w),\qquad
 v_\epsilon=y+\frac{\epsilon}{1-\epsilon}(y-w).
\]
Since $y\in\relint(Y)$ and $w\in\aff(Y)$, we have $v_\epsilon\in Y$ for
small enough $\epsilon$. Therefore
$y_\epsilon=(1-\epsilon)v_\epsilon+\epsilon q\in Q$.
Affineness gives $\sigma(y_\epsilon)=\sigma(y)=\bar z\in\Z^s$ and
$h(y_\epsilon)=\epsilon h(q)>0$, contradicting validity.
\end{proof}

\begin{lemma}[Keeping the slices separate]\label{lem:slicing}
Under the assumptions of Lemma~\ref{lem:latticefree}, if
$Q\cap H\cap\sigma^{-1}(\Z^s)$ is nonempty, it is the union of at most
$b(s)$ mixed-integer sets, each obtained from $Q$ by adding equations and
using $s-1$ integrality conditions. If $Q$ has a description with $m$
inequalities, each resulting polyhedron uses those same $m$ inequalities.
\end{lemma}

\begin{proof}
Let $K$ be the set in Lemma~\ref{lem:latticefree}. If it is full-dimensional,
then $\intr(\cl K)=\intr(K)$, so Theorem~\ref{thm:flatness} applies to
$\cl K$. Dividing the resulting integer direction by the greatest common
divisor of its coordinates makes it primitive and can only decrease its width.
There are at most $b(s)$ integer values of $v^Tz$ on $K$.

If $K$ is not full-dimensional, its affine hull is generated by integer points:
choose an affine basis from the integer set whose convex hull is $K$.
Thus there is a primitive $v\in\Z^s\setminus\{0\}$ and an integer $i_0$
such that $v^Tz=i_0$ on $K$. In this case only one slice is needed.

In either case, let $T\subseteq\Z$ be the finite set of attained integer
values, with $|T|\leq b(s)$. Extend $v^T$ to a unimodular matrix
$U\in\Z^{s\times s}$, and let $\sigma'$ consist of the last $s-1$
components of $U\sigma$. For $i\in T$, set
\[
 Q_i=Q\cap H\cap\{y:v^T\sigma(y)=i\}.
\]
Unimodularity gives the exact set identity
\begin{equation}\label{eq:exactunion}
 Q\cap H\cap\sigma^{-1}(\Z^s)
 =\bigcup_{i\in T}\bigl(Q_i\cap(\sigma')^{-1}(\Z^{s-1})\bigr).
\end{equation}
No convex hull is taken in this replacement. When $s=1$, the map $\sigma'$
has no components and imposes no integrality condition.
\end{proof}

\subsection{Disjunctive programming with bounded images}
We will use disjunctive programming only on linear relaxations whose images
are already contained in a bounded target polytope. The extension spaces
themselves need not be bounded.

\begin{lemma}\label{lem:disjunction}
For $a=1,\ldots,r$, let $Q_a$ be a nonempty polyhedron described by $m_a$
inequalities and arbitrary equations, and let $\pi_a$ be an affine map.
If each $\pi_a(Q_a)$ is bounded, then
\[
 R=\conv\Bigl(\bigcup_{a=1}^r\pi_a(Q_a)\Bigr)
\]
has a linear extended formulation with at most $\sum_a(m_a+1)$ inequalities.
\end{lemma}

\begin{proof}
This is the usual disjunctive formulation~\cite{Balas1998}; we verify the
bounded-image detail. Write
$Q_a=\{y:A_a y\leq d_a,\ D_a y=e_a\}$ and
$\pi_a(y)=T_a y+t_a$. Use variables $y^a,\lambda_a$ and impose
\begin{align*}
 A_a y^a&\leq d_a\lambda_a,& D_a y^a&=e_a\lambda_a,&\lambda_a&\geq0
 &&(a=1,\ldots,r),\\
 \sum_a\lambda_a&=1,& x&=\sum_a(T_a y^a+t_a\lambda_a).
\end{align*}
If $\lambda_a>0$, then $y^a/\lambda_a\in Q_a$. If $\lambda_a=0$,
then $y^a$ is a recession direction of $Q_a$. Boundedness of $\pi_a(Q_a)$
and nonemptiness of $Q_a$ imply $T_a y^a=0$. Thus every feasible $x$ is
a convex combination of points in the sets $\pi_a(Q_a)$. Conversely, every
such convex combination has a feasible lift, taking $y^a=0$ when its weight
is zero. The inequality count is $\sum_a(m_a+1)$.
\end{proof}

\section{An induction over families of formulations}\label{sec:families}

\subsection{The invariant and its potential}
At every stage we retain a vertex set $V$ and a finite family of triples
$(Q_a,\sigma_a,\pi_a)$, indexed by $a\in\mathcal A$, such that
$\sigma_a$ has $s_a$ components, $\pi_a$ takes values in $\R^{E(V)}$,
and
\begin{equation}\label{eq:invariant}
 \PM(V)=\conv\Bigl(\bigcup_{a\in\mathcal A} S_a\Bigr),\qquad
 S_a=\pi_a\bigl(Q_a\cap\sigma_a^{-1}(\Z^{s_a})\bigr).
\end{equation}
Each $Q_a$ has at most $m$ inequalities, together with equations. We discard
members with an empty mixed-integer set. In particular, each retained $Q_a$
is nonempty, and the family itself is nonempty because $\PM(V)$ is nonempty.
Equation~\eqref{eq:invariant} also implies $S_a\subseteq\PM(V)$ for every
member. No individual $S_a$ is required to generate all of $\PM(V)$.

Define integers
\begin{equation}\label{eq:B}
 B(0)=1,\qquad B(s)=1+b(s)B(s-1)\quad(s\geq1),
\end{equation}
and associate with a family the potential
\begin{equation}\label{eq:potential}
 \Psi=\sum_{a\in\mathcal A}B(s_a).
\end{equation}
The potential bounds the number of members, since $B(s)\geq1$. Replacing
a member with $s\geq1$ integrality conditions by at most $b(s)$ members
with $s-1$ conditions decreases $\Psi$ by at least one. Deleting a member
also decreases $\Psi$ by at least one.

\subsection{One step}

\begin{lemma}\label{lem:step}
Suppose~\eqref{eq:invariant} holds for a family of $r$ members, and let
$W\subseteq V$ satisfy $|W|\geq2$ and
\begin{equation}\label{eq:hardW}
 2^{c|W|}>r(m+1).
\end{equation}
There is a new family satisfying~\eqref{eq:invariant} on $V\setminus W$,
with at most $m$ inequalities per member and potential at most $\Psi-1$.
\end{lemma}

\begin{proof}
Let $\rho_a$ be $\pi_a$ followed by restriction to $E(W)$. If
$\rho_a(Q_a)\subseteq\PM(W)$ for every $a$, then
\[
 \PM(W)=\conv\Bigl(\bigcup_a\rho_a(Q_a)\Bigr).
\]
Indeed, containment from right to left follows from the assumption, and
containment from left to right follows by projecting~\eqref{eq:invariant}
onto $E(W)$. Every image is bounded. Lemma~\ref{lem:disjunction} would
therefore give an extended formulation of $\PM(W)$ of size at most
$r(m+1)$, contradicting Theorem~\ref{thm:rothvoss}
and~\eqref{eq:hardW}.

Consequently there is a member $a_*$ and a point $q\in Q_{a_*}$ whose
image violates a facet inequality $\alpha^Tx\leq\beta$ of $\PM(W)$.
The matching polytope on $W$ is full-dimensional in its edge space, so such
a facet exists. Let
\[
 h_a(y)=\alpha^T\rho_a(y)-\beta,\qquad H_a=\{y:h_a(y)=0\},
\]
and let
$F=\{x\in\PM(V):\alpha^Tx_{E(W)}=\beta\}$.
Each $h_a\leq0$ is valid on the mixed-integer set of member $a$, since
$S_a\subseteq\PM(V)$. Moreover,
\begin{equation}\label{eq:faceinvariant}
 F=\conv\Bigl(\bigcup_a
 \pi_a(Q_a\cap H_a\cap\sigma_a^{-1}(\Z^{s_a}))\Bigr).
\end{equation}
To see this, express a point of $F$ as a finite convex combination using
\eqref{eq:invariant}. Every point in that combination satisfies the supporting
inequality. Equality for the combination forces equality for every term
with positive weight. The reverse inclusion is immediate.

Restrict each $Q_a$ to $Q_a\cap H_a$. For $a\ne a_*$, retain its
integrality map. If the restricted mixed-integer set of $a_*$ is empty,
delete that member. Otherwise, $s_{a_*}\geq1$: a member with no integrality
conditions has every point of its relaxation in its mixed-integer set and
therefore cannot violate the valid inequality. Lemmas~\ref{lem:latticefree}
and~\ref{lem:slicing} replace the restricted member $a_*$ by at most
$b(s_{a_*})$ separate members with $s_{a_*}-1$ integrality conditions.
The exact union identity~\eqref{eq:exactunion} preserves
\eqref{eq:faceinvariant}.

Finally, compose every output map with projection onto $E(V\setminus W)$
and discard any empty mixed-integer members. Lemma~\ref{lem:matchingface}
proves the invariant for $\PM(V\setminus W)$. All restrictions added only
equations, so the inequality count of each member stays at most $m$.
The selected member either disappears or is replaced by children whose
total potential is at most $B(s_{a_*})-1$. Other members do not increase
the potential, proving the claimed decrease.
\end{proof}

\subsection{A quantitative bound}

\begin{theorem}\label{thm:quantitative}
Let $m,k\geq0$ be integers. If the matching polytope of $K_n$ has a MILEF
of size at most $(m,k)$, then
\begin{equation}\label{eq:quantitative}
 n<B(k)w(m,k),\qquad
 w(m,k)=\max\left\{2,\left\lfloor
 \frac{\log((m+1)B(k))}{c}\right\rfloor+1\right\},
\end{equation}
where $B$ is defined in~\eqref{eq:B} and $c$ is from
Theorem~\ref{thm:rothvoss}. In particular,
\begin{equation}\label{eq:asymptotictradeoff}
 n\leq 2^{O(k\log(k+1))}
 \bigl(1+\log(m+1)+k\log(k+1)\bigr),
\end{equation}
with absolute implicit constants.
\end{theorem}

\begin{proof}
Start with the single given formulation, so that the invariant holds and the
potential is $B(k)$. Suppose, to derive a contradiction, that
$n\geq B(k)w(m,k)$, and put $w=w(m,k)$.

After $t$ steps of Lemma~\ref{lem:step}, the graph has $n-tw$ vertices and
the potential is at most $B(k)-t$. In particular, the number $r$ of members
is at most $B(k)$. The definition of $w$ guarantees
\[
 2^{cw}>(m+1)B(k)\geq r(m+1).
\]
For every $t=0,\ldots,B(k)-1$, the assumed lower bound on $n$ leaves
at least $w$ vertices. Thus a set $W$ of size $w$ can be selected and
Lemma~\ref{lem:step} can be applied again. After $B(k)$ steps the
potential would be at most zero. This is impossible: the invariant describes
a nonempty matching polytope, so at least one member remains, and every
member contributes at least one to the potential. This proves
\eqref{eq:quantitative}. The argument also covers $k=0$; in that case a
violating member with no integrality conditions is already impossible.

Since $b(s)=O(s^{3/2})$, unfolding~\eqref{eq:B} gives
\[
 B(k)=1+\sum_{j=1}^k\prod_{t=j}^k b(t)
 \leq (k+1)\prod_{t=1}^k b(t)
 \leq (k+1)D^k(k!)^{3/2}
 =2^{O(k\log(k+1))}
\]
for an absolute constant $D$. Substituting this estimate into
\eqref{eq:quantitative} proves~\eqref{eq:asymptotictradeoff}.
\end{proof}

\begin{proof}[Proof of Theorem~\ref{thm:main}]
For sufficiently large $n$, the polynomial bound gives
$\log(m+1)=O_p(\log n)$. If $k\geq\log n$, the desired conclusion
already holds. Otherwise $k<\log n$, and~\eqref{eq:asymptotictradeoff}
implies, after taking logarithms,
\[
 \log n\leq O\bigl(k\log(k+1)\bigr)+O_p(\log\log n).
\]
For large enough $n$ this forces
$k\log(k+1)=\Omega_p(\log n)$. Using
$\log(k+1)=O(\log\log n)$ proves the result.
\end{proof}

\section{What the corrected argument establishes}\label{sec:scope}

The proof preserves the two geometric steps of the earlier approach: a
violated inequality produces a lattice-free set on its equality hyperplane,
and a face supported on a small vertex set leaves a smaller matching polytope
as a projection. The correction concerns how these steps are composed.
An individual slice need not describe a matching polytope. The family invariant
allows this, while the common face restriction preserves the target for the
family as a whole.

Disjunctive programming enters only in the test that finds a violating
member. Under the contrary assumption, the images of all linear relaxations
already lie in the target matching polytope, and their convex hull can safely
be used. Integrality conditions are not imposed after this convexification.
The proof terminates by a potential contradiction; it does not require
combining mixed-integer branches into a single MILEF with fewer integer
variables.

The weaker bound reflects the number of branches. Along a lineage at most
$k$ integrality conditions are removed, but there may be
$2^{O(k\log(k+1))}$ nodes to process. Our proof charges a separate loss of
graph vertices whenever a member is processed. It does not show that one
face makes progress in all members simultaneously, or that these losses can
be charged only to the depth of the branching tree. Accordingly, it does not
recover the $\Omega(\sqrt{n/\log n})$ bound claimed in the earlier version.

The stronger lower bounds of~\cite{CWZ2018} remain the relevant quantitative
results. This revision supplies a corrected account of the original method
and an alternative, more specialized proof of a weaker consequence.

\section*{Acknowledgements}
We thank Stefan Weltge and Gennadiy Averkov for the discussions on flatness
theorems acknowledged in the earlier version. We also acknowledge the
comments of the reviewers of that version.

\appendix
\section{The error in the earlier proof}\label{app:error}

\subsection{Location and the invalid equality}
In arXiv:1611.00707v1~\cite{HWZ2016}, the proof of Lemma~8 on printed
page~11 contains a chain of equalities intended to justify eliminating one
integer variable. After correcting the incidental indexing of the slices,
the problematic step has the following form. Let
\[
 S=\bigcup_{i=\ell}^{u}\bigl(P\cap\{x:x_\tau=i\}\bigr),
 \qquad J'=J\setminus\{\tau\},
\]
and write $L_{J'}$ for the points with integer coordinates in $J'$. The proof
uses
\begin{equation}\label{eq:false}
 \proj_I\conv(S\cap L_{J'})
 \stackrel{\text{claimed}}{=}
 \proj_I\conv\bigl(\conv(S)\cap L_{J'}\bigr).
\end{equation}
Only the left-to-right inclusion is automatic. Convex combinations of points
that are fractional in the remaining integer coordinates may have integer
averages. Such points can enter the right-hand side without being convex
combinations of feasible mixed-integer points of the slices. Projection does
not in general remove the discrepancy.

\subsection{A bounded rational counterexample}
Consider the polytope
\[
 P=\left\{(x,t,z)\in\R^3:
 0\leq x\leq1,\quad0\leq t\leq1,\quad z=t+\frac{x}{2}\right\}.
\]
The output is $x$, and $t,z$ are required to be integer. Since
$x/2=z-t$ is an integer in $[0,1/2]$, every mixed-integer feasible point
has $x=0$. Hence
\begin{equation}\label{eq:counterleft}
 \proj_x\conv(P\cap\{t,z\in\Z\})=\{0\}.
\end{equation}
The variable $t$ takes only the values $0$ and $1$. Its two slices contain,
respectively,
\[
 a=(1,0,1/2),\qquad b=(1,1,3/2).
\]
Their midpoint is $(1,1/2,1)$, which satisfies the remaining integrality
condition $z\in\Z$ and projects to $x=1$. In fact, for every $x\in[0,1]$
the point $(x,1-x/2,1)$ belongs to the convex hull of the two slices.
Consequently, for $S=(P\cap\{t=0\})\cup(P\cap\{t=1\})$,
\begin{equation}\label{eq:counterright}
 \proj_x\conv\bigl(\conv(S)\cap\{z\in\Z\}\bigr)=[0,1].
\end{equation}
Equations~\eqref{eq:counterleft} and~\eqref{eq:counterright} contradict
\eqref{eq:false}. The example has a bounded rational relaxation, a primitive
elimination direction, and nonempty mixed-integer sets in both slices.
Thus these additional properties do not validate the construction.

This counterexample refutes the construction and the equality used in the
proof. It does not, by itself, refute the existential statement of the old
Lemma~8: in the example the true projected polytope $\{0\}$ has a trivial
linear formulation. The present revision does not use that lemma.

\subsection{The valid replacement and its consequences}
For $P_i=P\cap\{x_\tau=i\}$, the valid identity is
\[
 \conv(P\cap L_J)
 =\conv\left(\bigcup_{i=\ell}^{u}
       \conv(P_i\cap L_{J'})\right),
\]
provided the slices include all values attained by the eliminated integer
coordinate. This identity retains integrality within each slice before any
convexification across slices. It does not bound the linear extension
complexity of the mixed-integer hull of an individual slice.

Lemma~\ref{lem:slicing} instead retains the union as a family. The proof
of Lemma~\ref{lem:step} then applies a common face restriction to that
family, and the potential~\eqref{eq:potential} accounts for the extra
members. This is the replacement for the old Lemma~8 and its use in the
old Lemma~5. The lower bound and its algorithmic interpretation have been
weakened throughout the manuscript to match this accounting.

There is a separate hypothesis mismatch in the old application of Lemma~8:
flatness bounds the chosen direction on the mixed-integer hull of the face,
whereas that lemma assumes the bound on the full relaxation. Intersecting the
relaxation with the corresponding two slab inequalities would fix this
mismatch without removing any mixed-integer feasible points, but would not
fix~\eqref{eq:false}. Our slicing lemma needs bounds only on the
mixed-integer points and adds slice equations directly, so it avoids that
additional issue.

% The bibliography below is embedded so this file is standalone.
% Entries retained from main-online.bbl are followed by the two revision references.
\begin{thebibliography}{99}
\bibitem{bader_2016_mixed}
J.~Bader, R.~Hildebrand, R.~Weismantel, and R.~Zenklusen.
\newblock Mixed integer reformulations of integer programs and the affine
{TU}-dimension of a matrix, 2016.
\newblock \url{https://arxiv.org/abs/1508.02940}.

\bibitem{Balas1998}
E.~Balas.
\newblock Disjunctive programming: Properties of the convex hull of feasible points.
\newblock \emph{Discrete Applied Mathematics}, 89(1--3):3--44, 1998.

\bibitem{Banaszczyk1999}
W.~Banaszczyk, A.~E. Litvak, A.~Pajor, and S.~J. Szarek.
\newblock The flatness theorem for nonsymmetric convex bodies via the local
theory of Banach spaces.
\newblock \emph{Mathematics of Operations Research}, 24(3):728--750, 1999.

\bibitem{CWZ2018}
A.~Cevallos, S.~Weltge, and R.~Zenklusen.
\newblock Lifting linear extension complexity bounds to the mixed-integer setting.
\newblock In \emph{Proceedings of the Twenty-Ninth Annual ACM--SIAM Symposium
on Discrete Algorithms (SODA)}, pages 788--807, 2018.
\newblock \href{https://doi.org/10.1137/1.9781611975031.51}{doi:10.1137/1.9781611975031.51}.

\bibitem{HWZ2016}
R.~Hildebrand, R.~Weismantel, and R.~Zenklusen.
\newblock Extension complexity lower bounds for mixed-integer extended formulations.
\newblock arXiv:1611.00707v1, 2016. Withdrawn in 2022.
\newblock \url{https://arxiv.org/abs/1611.00707v1}.

\bibitem{kaibel_2015_lower}
V.~Kaibel and S.~Weltge.
\newblock Lower bounds on the sizes of integer programs without additional variables.
\newblock \emph{Mathematical Programming}, 154(1):407--425, 2015.

\bibitem{Khinchin1948}
A.~Khinchin.
\newblock A quantitative formulation of Kronecker's theory of approximation
(in Russian).
\newblock \emph{Izvestiya Akademii Nauk SSR Seriya Matematika}, 12:113--122, 1948.

\bibitem{lenstra_integer_1983}
H.~W. Lenstra, Jr.
\newblock Integer programming with a fixed number of variables.
\newblock \emph{Mathematics of Operations Research}, 8:538--548, 1983.

\bibitem{miller_1960_integer}
C.~E. Miller, A.~W. Tucker, and R.~A. Zemlin.
\newblock Integer programming formulation of traveling salesman problems.
\newblock \emph{Journal of the ACM}, 7(4):326--329, 1960.

\bibitem{oncan_2009_comparative}
T.~{\"O}ncan, I.~K. Alt{\i}nel, and G.~Laporte.
\newblock A comparative analysis of several asymmetric traveling salesman
problem formulations.
\newblock \emph{Computers \& Operations Research}, 36(3):637--654, 2009.

\bibitem{padberg_1991_analytical}
M.~Padberg and T.-Y. Sung.
\newblock An analytical comparison of different formulations of the travelling
salesman problem.
\newblock \emph{Mathematical Programming}, 52(1):315--357, 1991.

\bibitem{rothvoss_2014_matching}
T.~Rothvo{\ss}.
\newblock The matching polytope has exponential extension complexity.
\newblock In \emph{Proceedings of the 46th Annual ACM Symposium on Theory
of Computing (STOC)}, pages 263--272, 2014.

\bibitem{yannakakis_1991_expressing}
M.~Yannakakis.
\newblock Expressing combinatorial optimization problems by linear programs.
\newblock \emph{Journal of Computer and System Sciences}, 43(3):441--466, 1991.
\end{thebibliography}
\end{document}
